\documentclass{article}

\usepackage{fullpage}
\usepackage{url}
\usepackage{bbm}
\usepackage{graphicx}
\usepackage{amsmath}
\usepackage{physics}
\usepackage{mathtools}
\usepackage{bbold}
\usepackage{slashed}
\usepackage{pxfonts}
\usepackage{multicol}
\usepackage{mathptmx}
\usepackage{simplewick}
\usepackage{placeins}

\usepackage{lmodern}
\usepackage[T1]{fontenc}
\usepackage[fontsize=13.5pt]{fontsize}

\DeclareMathOperator{\erfi}{erfi}
\newcommand{\R}{\ensuremath{\mathbb{R}}}
\newcommand{\C}{\ensuremath{\mathbb{C}}}

\title{Quantum Field Theory I Excercise 5}
\author{Idan Stark}
\date{\today}

\begin{document}
\maketitle

\section{A theory of two scalar fields}
In this excercise we explore the theory with two real scalar fields $\phi_1, \phi_2$, in $3+1$ dimensions, with the lagrangian density:
\begin{equation}
    \mathcal{L} = \sum_{i=1,2}\left(
        \frac{1}{2}\left(\partial_\mu\phi_i\right)^2 - \frac{1}{2}m_i^2\phi_i^2
    \right) - \frac{g}{4}\phi_1^2\phi_2^2
\end{equation}
\subsection{Symmetries}
Assuming the masses are different, both fields have independent $\phi_i \rightarrow -\phi_i$ symmetries. Assuming they are the same, the lagrangian density becomes:
\begin{equation}
    \mathcal{L} = \sum_{i=1,2}\left(
        \frac{1}{2}\left(\partial_\mu\phi_i\right)^2
    \right) - \frac{1}{2}m^2(\phi_1^2 + \phi_2^2) - \frac{g}{4}\phi_1^2\phi_2^2
\end{equation}
Taking the vector $\vec{\phi} = \phi_i$ as a 2D vector in a fields space, we can see that the mass term is basically the squared magnitude of the vector and the interaction term is the squared area created by its coordinates. So, any linear transformation that preserves vector magnitude and area will be a symmetry. To preserve the magnitude, let us look at rotations in that fields space and require they preserve the interaction term:
\begin{equation}
    \begin{pmatrix}
        \phi_1 \\ \phi_2
    \end{pmatrix} \rightarrow
    \begin{pmatrix}
        \phi_1' \\ \phi_2'
    \end{pmatrix} = \begin{pmatrix}
        \cos \alpha & \sin \alpha \\ -\sin\alpha & \cos \alpha
    \end{pmatrix}
    \begin{pmatrix}
        \phi_1 \\ \phi_2
    \end{pmatrix}
\end{equation}
First, let us prove it preserves the first two terms. This is simple. First, we notice:
\begin{equation}
\begin{gathered}
    \phi_1'^2 = (\cos\alpha \phi_1 + \sin\alpha \phi_2)^2 = \cos^2\alpha \phi_1^2 + 2\cos\alpha\sin\alpha\phi_1\phi_2 + \sin^2\alpha\phi_2^2
    \\
    \phi_1'^2 = (\cos\alpha \phi_2 - \sin\alpha \phi_1)^2 = \cos^2\alpha \phi_2^2 - 2\cos\alpha\sin\alpha\phi_1\phi_2 + \sin^2\alpha\phi_1^2
\end{gathered}
\end{equation}
Which means
\begin{equation}
    \phi_1'^2 + \phi_2'^2 = (\cos^2\alpha + \sin^2\alpha)\phi_1^2 + (\cos^2\alpha + \sin^2\alpha)\phi_2^2 = \phi_1^2 + \phi_2^2
\end{equation}
And since $\partial_\mu \alpha = 0$:
\begin{equation}
\begin{gathered}
    \left(\partial_\mu\phi_1'\right)^2 + \left(\partial_\mu\phi_2'\right)^2 = \\ = \cos^2\alpha \left(\partial_\mu\phi_1\right)^2 + 2\cos\alpha\sin\alpha\partial_\mu\phi_1\partial^\mu\phi_2 + \sin^2\alpha\left(\partial_\mu\phi_2\right)^2 + \\ +
    \cos^2\alpha \left(\partial_\mu\phi_2\right)^2 - 2\cos\alpha\sin\alpha\partial_\mu\phi_1\partial^\mu\phi_2 + \sin^2\alpha\left(\partial_\mu\phi_1\right)^2 = \\ =
    \left(\partial_\mu\phi_1\right)^2 + \left(\partial_\mu\phi_2\right)^2
\end{gathered}
\end{equation}
Now, let us require that the interaction term is preserved:
\begin{equation}
\begin{gathered}
    \phi_1'\phi_2' = (\cos\alpha\phi_1 + \sin\alpha\phi_2)(\cos\alpha\phi_2 - \sin\alpha\phi_1) = \\ = (\cos^2\alpha - \sin^2\alpha)\phi_1\phi_2 + \cos\alpha\sin\alpha(\phi_2^2 - \phi_1^2) = \\ =
    \cos2\alpha\phi_1\phi_2 + \frac{1}{2}\sin2\alpha(\phi_2^2 - \phi_1^2) = \phi_1\phi_2
\end{gathered}
\end{equation}
This means:
\begin{equation}
    \cos2\alpha = 1 \qquad \sin2\alpha = 0
\end{equation}
\begin{equation}
    \alpha = \frac{\pi n}{2}
\end{equation}
So, the additional symmetries of the system are any 90 degrees rotation in the fields space.

\subsection{Effects of the symmetries}
The global symmetry that applies even if the masses are different can be written as:
\begin{equation}
    S[\phi_1, \phi_2] = S[-\phi_1, \phi_2] = S[\phi_1, -\phi_2] = S[-\phi_1, -\phi_2]
\end{equation}
It's easy to see then, under this symmetry transformation, the in the path integral the action will remain unchanged while the field integration measure will flip sign. This directly means that:
\begin{equation}
\begin{gathered}
    \langle \phi_1 \rangle =
    \int \mathcal{D}\phi_1\mathcal{D}\phi_2 \phi_1 e^{iS[\phi_1, \phi_2]} = \\ =
    \int \mathcal{D}\phi_1\mathcal{D}\phi_2 \phi_1 e^{iS[-\phi_1, \phi_2]} =
    \int \mathcal{D}\phi_1\mathcal{D}\phi_2 (-\phi_1) e^{iS[\phi_1, \phi_2]} =
    -\langle \phi_1 \rangle
\end{gathered}
\end{equation}
And similarly for $\phi_2$. So:
\begin{equation}
    \langle \phi_i \rangle = - \langle \phi_i \rangle = 0
\end{equation}
Similarly:
\begin{equation}
\begin{gathered}
    \langle \phi_1 \phi_2 \rangle =
    \int \mathcal{D}\phi_1\mathcal{D}\phi_2 \phi_1 \phi_2 e^{iS[\phi_1, \phi_2]} = \\ =
    \int \mathcal{D}\phi_1\mathcal{D}\phi_2 \phi_1 \phi_2 e^{iS[-\phi_1, \phi_2]} =
    \int \mathcal{D}\phi_1\mathcal{D}\phi_2 (-\phi_1) \phi_2 e^{iS[\phi_1, \phi_2]} =
    -\langle \phi_1 \phi_2 \rangle
\end{gathered}
\end{equation}
And so:
\begin{equation}
    \langle \phi_1(x) \phi_2(y) \rangle = - \langle \phi_1(x) \phi_2(y) \rangle = 0
\end{equation}
\subsection{Renormalizable theory}
Following the original symmetry, we can only get items where each field appears with an even power:
\begin{equation}
    \mathcal{L} = \mathcal{L}_0 + A_1\phi_1^4 + A_2\phi_2^4 + B_1\phi_1^6 + \cdots + a_1 \left(\partial_\mu \phi_1\right)^4 + a_2 \left(\partial_\mu \phi_2\right)^4 + a_i \left(\partial_\mu \phi_1\right)^2 \left(\partial_\mu \phi_2\right)^2 + \cdots
\end{equation}
Now, every term with mass dimension greater then four will have a coupling paramter that must scale with a negative power of the cutoff $\Lambda$, so the only remaining terms are those with total mass dimension 4 or less. This leaves out any combination of the derivatives, since a squared derivative of the field already has mass dimension 4, and of course and powers of the fields greater then 4, leaving only two additional terms:
\begin{equation}
    \mathcal{L} = \sum_{i=1,2}\left(
        \frac{1}{2}\left(\partial_\mu\phi_i\right)^2 - \frac{1}{2}m_i^2\phi_i^2
        + \frac{\lambda_i}{4}\phi_i^4
    \right) - \frac{g}{4}\phi_1^2\phi_2^2
\end{equation}
\subsection{Terms when $m_1 = m_2$}
When the masses are the same:
\begin{equation}
    \mathcal{L} = \sum_{i=1,2}\left(
        \frac{1}{2}\left(\partial_\mu\phi_i\right)^2 - \frac{1}{2}m^2\phi_i^2
        + \frac{\lambda_i}{4}\phi_i^4
    \right) - \frac{g}{4}\phi_1^2\phi_2^2
\end{equation}
This is just a renormalizaton of the original theory, so we impose the same symmetries. To follow them, we must force:
\begin{equation}
    S[\phi_1, \phi_2] = S[-\phi_2, \phi_1] = S[\phi_2, \phi_1]
\end{equation}
This directly forces $\lambda_1 = \lambda_2 = \lambda$, so the action now has one additional coupling, in comparison to the two there were in the different theories masses:
\begin{equation}
    \mathcal{L} = \sum_{i=1,2}\left(
        \frac{1}{2}\left(\partial_\mu\phi_i\right)^2 - \frac{1}{2}m^2\phi_i^2
        + \frac{\lambda}{4}\phi_i^4
    \right) - \frac{g}{4}\phi_1^2\phi_2^2
\end{equation}
\subsection{Imaginary part of $\Sigma(p^2)$}
We remember (thanks to the hint) that an imaginary part of the 1PI corection $\Sigma(p^2)$ comes from physical intermidate states. We expect that intermidate states can only occour if center of mass energy is sufficient to create additional particles.  In the interaction from the original theory, the only way to get physical intermidiate states with minimal vertices, is through the $\theta$ diagram, creating three intermidiate particles.
% TODO draw the theta diagram, example using 2m_1 and m_2
\newline
If the masses are equal $m_1 = m_2 = m$, the the total energy required is the mass of the three particles:
\begin{equation}
    p^2 \ge 9m^2
\end{equation}
Similarly, if $m_1^2 < m_2^2$, the above diagram requires $2m_1 + m_2$, and will contribute only if the three particles, two with mass $m_1$ and one with mass $m_1$ can be created:
\begin{equation}
    p^2 \ge (2m_1 + m_2)^2
\end{equation}
\subsection{Possible processes}
Since this is a $d = 3 + 1$ theory, we only look at processes for $N = 1, 2, 3, 4$, but as we saw in the second section, all $N = 1, 3$ processes vanish, leaving us with $N = 2, 4$. Additionally, we saw that the 2-point function mixing the fields also vanishes, leaving only the 2-point functions to be $\langle \phi_1(x) \phi_1(y) \rangle$ and $\langle \phi_2(x) \phi_2(y) \rangle$. These processes will have superficial DOF:
\begin{equation}
    D = d + \left(n\left(\frac{d - 2}{2}\right) - d\right)V - \frac{d - 2}{2}N = 
    4 + \left(4 - 4\right)V - 2 = 2
\end{equation}
And so these 2-point functions diverge quadratically.
\newline
Finally, for the $N = 4$ processes the superficial DOF is:
\begin{equation}
    D = d + \left(n\left(\frac{d - 2}{2}\right) - d\right)V - \frac{d - 2}{2}N = 
    4 + \left(4 - 4\right)V - 4 = 0
\end{equation}
And so all 4-point correlation functions diverge logarithmically.
\newline
For the renormalized perturbation theory, we will need to use five counter terms: one for each of the masses ($\delta m_1, \delta m_2$) and one for the each of the interaction ($\delta g, \delta \lambda_1, \delta \lambda_2$).
\subsection{Renormalization of $g$}
The original Lagrangian is written using the bare terms. Taking $\phi_1 \rightarrow \sqrt{Z_1}\phi_{p1}$ and $\phi_2 \rightarrow \sqrt{Z_2}\phi_{p2}$:
\begin{equation}
    \mathcal{L} = \sum_{i=1,2}Z_i\left(
        \frac{1}{2}\left(\partial_\mu\phi_{pi}\right)^2 - \frac{1}{2}m_i^2\phi_{pi}^2
    \right) - \frac{g}{4}Z_1Z_2\phi_{p1}^2\phi_{p2}^2
\end{equation}
Comparing this to the physical Lagrangian, defined using the scattering amplitude at zero spatial momentum, we get the additional counter terms:
\begin{equation}
\begin{gathered}
    \mathcal{L} = \sum_{i=1,2}\left(
        \frac{1}{2}\left(\partial_\mu\phi_{pi}\right)^2 - \frac{1}{2}m_{pi}^2\phi_{pi}^2
    \right) - \frac{g_p}{4}\phi_{p1}^2\phi_{p2}^2  + \\ + \sum_{i=1,2}\left(
        \frac{1}{2}\delta_{Z_i}\left(\partial_\mu\phi_{pi}\right)^2 - \frac{1}{2}\delta_{m_i}^2\phi_{pi}^2
    \right) - \frac{\delta_g}{4}\phi_{p1}^2\phi_{p2}^2
\end{gathered}
\end{equation}
Where:
\begin{equation}
    \delta_{Z_i} = Z - 1
    \qquad
    \delta_{m_i} = m_i^2 Z - m_{pi}^2
    \qquad
    \delta_g = gZ - g_p
\end{equation}
Now, let us look at the resulting diagrams. For the renormalization of $g$, we only need to look at $\delta_g$, which means that there is only one additional vertex, and it's identical to the original one.
\newline
In the scattering of two particles of different types $\phi_1\phi_2 \rightarrow \phi_1\phi_2$ these vertices create the following channels:
\begin{itemize}
    \item Two channels $s$ and $t$, for each vertex type, with orders $g^2\delta_g \sim g^2$ and $\delta_g^2 \sim g^4$ which can be ignored, based on the vertex. Each of the channels has two diagrams matching it, based on which internal propegator is $\phi_1$ and which is $\phi_2$.
    \item Six mixed channels $s, t, u$, where one vertex is $g_p$ and the other is $\delta_g$, all of which are of order $g\delta_g \sim g^3$, and therefore ignored.
\end{itemize}
% TODO add diagrams
Each of these original diagrams gives the contribution $V(p^2)$, so the total contribution is:
\begin{equation}
    i\mathcal{M} = -ig_p + 2(-ig_p)^2\left(iV(s) + iV(t)\right) -i\delta_g
\end{equation}
This has the same structure as the $\phi^4$ theory we saw in class. To find $\delta_g$, we will use the renormalization condition defining the physical constant $g_p$:
\begin{equation}
    \mathcal{M}_p(p_i = (m, \vec{0})) = -g_p
\end{equation}
\begin{equation}
    -g_p + 2(-ig_p)^2\left(V(4m^2) + V(0)\right) -\delta_g = -g_p
\end{equation}
Using $g_p = g + \mathcal{O}(g^2)$:
\begin{equation}
    \delta_g = - 2g^2\left(V(4m^2) + V(0)\right)
\end{equation}
Which gives, for the energy scale $\mu$:
\begin{equation}
    \delta_g = - \frac{g^2}{16\pi^2}\int_0^1 dx \left(
        \frac{6}{\epsilon} - 3\gamma - \log\frac{m^2 - x(1 - x)4m^2}{4\pi\mu^2} - \log\frac{m^2}{4\pi\mu^2} + \mathcal{O}(\epsilon)\right)
\end{equation}
\subsection{The beta function}
Defining the physical coupling term as instructed:
\begin{equation}
    g_p = -\mathcal{M}(s = t = u = M^2) = g + 4g^2V(M^2) + \delta_g
\end{equation}
So, the beta function is:
\begin{equation}
    \beta = M^2 \frac{dg_p}{dM^2} = M^24g^2V'(M^2)
\end{equation}
As we saw in class:
\begin{equation}
    V'(M^2) = \frac{1}{16\pi^2M^2} + \mathcal{O}\left(
        g_p^3, \frac{m^2}{M^2}
    \right)
\end{equation}
And so:
\begin{equation}
    \beta = \frac{g^2}{4\pi^2} + \mathcal{O}(g^3)
\end{equation}
\subsection{The beta function for the $\lambda_i$ couplings}
Looking at the $\phi_1 \phi_1 \rightarrow \phi_1 \phi_1$ and $\phi_2 \phi_2 \rightarrow \phi_2 \phi_2$ correlation functions, they get a contribution from the $g\phi_1^2\phi_2^2$ interaction through the three $s, t, u$ channels at second order of $g^2$ (this time there is a $u$ channel). In the effective field theory dscribed in section 1.3, these are the $\lambda_i\phi_i^4$ vertices. 
% TODO add diagrams
Since this result comes from all three channels, the $\beta$ function will be identical to the one we got for $\lambda_p$ in the $\phi^4$ theory:
\begin{equation}
    \beta = \frac{3}{16\pi^2}g^2 + \mathcal{O}(g^3)
\end{equation}
\subsection{Special symmetry}
In the equal masses case from section 1.4:
\begin{equation}
    \mathcal{L} = \sum_{i=1,2}\left(
        \frac{1}{2}\left(\partial_\mu\phi_i\right)^2 - \frac{1}{2}m^2\phi_i^2
        + \frac{\lambda}{4}\phi_i^4
    \right) - \frac{g}{4}\phi_1^2\phi_2^2
\end{equation}
There is a special case when $\lambda = -\frac{g}{2}$, then, we can write:
\begin{equation}
\begin{gathered}
    \mathcal{L} = \sum_{i=1,2}\left(
        \frac{1}{2}\left(\partial_\mu\phi_i\right)^2 - \frac{1}{2}m^2\phi_i^2
    \right) - \frac{g}{8}\left[\phi_1^4 + 2\phi_1^2\phi_2^2 + \phi_2^4\right] = \\ =
    \sum_{i=1,2}\left(
        \frac{1}{2}\left(\partial_\mu\phi_i\right)^2 - \frac{1}{2}m^2\phi_i^2
    \right) - \frac{g}{8}\left[\phi_1^2 + \phi_2^2\right]^2
\end{gathered}
\end{equation}
In this case, we can see that any rotation in the fields space is a symmetry, i.e. the system has $U(2)$ symmetry.
\newline

\newline
Since we have in this effective theory $(\phi_1^2 + \phi_1^2)$ and $(\phi_1^2 + \phi_1^2)^2$, it is completely possible that in low energy cutoff we might get that $(\phi_1^2 + \phi_1^2)^3$ is non negligable as well.
\end{document}
